
\begin{abstract}

% 摘要只保留有信息量的句子。完成后在草稿中逐句标注
% [任务] [模型] [算法] [结果] [边界]，定稿时删除标注。

% ==== 研究对象与全文主线（1--2 句）====
本文针对……问题，研究……。全文以……为统一建模主线。

% ==== 逐问摘要：通常每问 2--3 个短句，不写成长并列句 ====
\textbf{对于问题一，}题目要求……，核心困难在于……。
为此建立……模型，并采用……方法求解。
结果得到……（关键数值或明确结论），经……验证……。

\textbf{对于问题二，}题目要求……。
建立……模型，利用……算法处理……。
最终……（关键数值），其结论范围为……。

\textbf{对于问题三，}题目要求……。
针对……难点，构建……，并通过……求解。
结果表明……（关键数值或可核查结论）。

\textbf{对于问题四，}题目要求……。
将问题转化为……模型，采用……算法获得……。
结合……下界或检验，可得……。

% ==== 结尾仅总结已被正文支持的统一贡献；无必要可删除 ====
综上，本文……。其中严格最优、启发式上界、预测结果或情景结论应明确区分。

\keywords{关键词1；关键词2；关键词3；关键词4}

\end{abstract}
